\documentclass[10pt]{article}
\usepackage[margin=20mm]{geometry}
\usepackage{amsmath,amssymb,fontspec,unicode-math,xurl,hyperref}
\IfFontExistsTF{DejaVu Serif}{\setmainfont{DejaVu Serif}}{\setmainfont{Latin Modern Roman}}
\setmathfont{latinmodern-math.otf}
\setlength{\parindent}{0pt}\setlength{\parskip}{6pt}\sloppy
\begin{document}
Date:\par
Title: Scientific verification report
I25-40131 / 0028\par
Position: Contextual Integrity is Inadequately Applied to Language Models\par
Yan Shvartzshnaider; Vasisht Duddu\par
ICML 2025 / PMLR 267, pp. 82200–82210; publisher PDF, 11 pages\par
Minor — one verified overstatement of prior experimental coverage\par
The statement that none of three cited studies accounted for prompt variation conflicts with a pre-existing paraphrasing experiment in one of those studies. This is a bounded literature-characterization error; the paper’s central distinction between contextual integrity (CI) and narrower privacy proxies remains supported within the reviewed scope.\par
\# 1. Identity, version and review scope\par
The original publisher PDF was retained unchanged. Its title, authors and proceedings identity match the task and acquisition receipt. SHA-256: `32a0ea673672ac12848e878c8b94d06a7f1500daedfad9cf819fc41e8ff81589`. Page references below use one-based PDF pages. Main text and Tables 1–2 were read; references occupy pp. 9–11 and there is no included appendix. Theoretical claims and the potentially consequential blanket statement in Section 4 received targeted primary-source checks.\par
The inspected arXiv record is 2501.19173, with v1 dated 31 January 2025 and v2 dated 15 May 2025. This review applies to the hashed proceedings PDF; identity with either arXiv PDF was not assumed. No correction was identified in the inspected publisher/arXiv records. The OpenReview response was a browser challenge, so reviews and revision discussions were unavailable.\par
\# 2. Problem 1: Prior paraphrasing experiment omitted\par
F28-01 — Minor; high confidence. Section 4, p. 7, under “Paraphrasing Prompt Sensitivity,” states that none of Mireshghallah et al. (2023), Ghalebikesabi et al. (2024), or Shao et al. (2024) accounted for prompt variation. The cited Ghalebikesabi paper provides a direct counterexample in its historical arXiv:2408.02373v1 version, dated 5 August 2024, before the reviewed paper.\par
In Operationalizing Contextual Integrity in Privacy-Conscious Assistants, Section 4.2 (PDF p. 6) varies form titles, descriptions and field wording. Section 5 (pp. 8–9) analyzes robustness across paraphrased forms; Figure 5 (p. 9) plots within-form-group variation in privacy leakage and utility. Appendix C.2 (pp. 33–34) gives paraphrasing instructions. These forms enter the assistant’s decision input. The page-9 figure and caption were visually inspected, and the exact historical PDF was archived with SHA-256 `6a6764795f9d19ac2cb7b87cc4655388c5d7550fc7eb4ee317a85bd009320c0c`.\par
The strongest alternative reading is that the reviewed authors meant variation only in a fixed instruction or assessment template, excluding the wording of form inputs. Section 4 does not state that restriction. The finding therefore concerns the unrestricted published assertion, not a claim that all possible prompt dimensions were previously tested. It is not a missing-citation complaint: the counterexample is already cited.\par
Supported impact and repair: replace the blanket statement with an explicit acknowledgement that Ghalebikesabi et al. tested form paraphrasing, then identify any distinct instruction-template variation left untested. The assertion of uniformly absent paraphrasing controls must be narrowed. This does not establish adequate coverage of position bias, validate every CI-tenet classification, reproduce the prior experiment, or overturn the social-theoretical position. No Major issue is supported.\par
\# 3. Core argument and candidate objections\par
Claim / location | Check and result\par
Four CI tenets; Section 2, pp. 2–3 | Consistent with Nissenbaum’s 2019 primary account: appropriate flows, contextual norms, parameterized flows, and ethical assessment of departures. No Issue Found in this conceptual correspondence.\par
Legal rules as privacy norms; Section 3 and Table 2, p. 8 | The heading alone could appear absolute. Table 2 and Section 5 explicitly allow institutions that reflect contextual values. A claim that the paper prohibits all legal embodiments of norms is not supported.\par
Partial use of CI; Section 5, p. 8 | The paper acknowledges useful applications of parts of the framework with limited claims. This qualification avoids interpreting all partial implementations as intrinsically worthless.\par
ConfAIde flow representation; Section 3 / Table 1 | The targeted 2024 primary version acknowledges five CI parameters but uses a three-factor Tier-2 template. The reviewed critique has a concrete target; no additional verified error was established.\par
Experimental hygiene; Section 4, p. 7 | Separating repeated-prompt, paraphrasing and option-position variation is useful. F28-01 affects the literature summary for the second dimension; the other universal negatives were not exhaustively verified.\par
Nissenbaum, Contextual Integrity Up and Down the Data Food Chain (2019), pp. 224, 227 and 231–234, supports the distinction between descriptive departures from prevailing norms and ethically illegitimate flows. It also permits partial use of the framework. The reviewed paper’s qualified discussion is compatible with this account. This is a conceptual consistency check, not an empirical test of privacy protection or a legal compliance opinion.\par
\# 4. Coverage, limits and reproducibility\par
Coverage includes the complete reviewed PDF and targeted checks of two cited LLM studies plus the primary CI account. It does not independently audit all nine rows and all four tenet labels in Table 1. Claims about every prior study’s lack of position-bias controls therefore remain Needs Evidence for an exhaustive verification. No new empirical experiment, equation-based performance claim, or author implementation is supplied by this position paper; no executable author code was required to test F28-01, and none was run.\par
The evidence bundle preserves publisher and arXiv metadata, primary-source PDFs/HTML, acquisition hashes, extracted text and page images. is reviewer-authored and only extracts, checks and renders public documents. Its successful check establishes the presence and historical date of the counterexample, not the correctness of the reported experimental measurements. OpenReview discussion and a complete all-study robustness audit remain unavailable or unperformed.\par
\# 5. Conclusion and repair status\par
Judgment / scope | Repair verification | Adoption / human status | Remaining work\par
F28-01: Minor | Acknowledging the existing form-paraphrase experiment is source-supported | Not requested / pending | Do not claim remaining prompt dimensions were exhaustively tested\par
Core CI theory: No Issue Found within scope | No repair proposed; qualifications preserved | Not requested / pending | Complete nine-study classification remains unaudited\par
\# Appendix. Sources, commands, evidence and roles\par
- Reviewed publisher record and PDF: https://proceedings.mlr.press/v267/shvartzshnaider25a.html ; original at .\par
- Reviewed-paper history: https://arxiv.org/abs/2501.19173 ; archived under sources/0028/.\par
- Ghalebikesabi et al. (2024), historical PDF: https://arxiv.org/pdf/2408.02373v1 ; Sections 4.2 and 5, Figure 5, Appendix C.2.\par
- Nissenbaum (2019), Theoretical Inquiries in Law 20(1):221–256, doi:10.1515/til-2019-0008 ; primary Cornell-hosted PDF archived as .\par
- Mireshghallah et al., Can LLMs Keep a Secret?, arXiv:2310.17884v2 (28 June 2024), Sections 2 and 3.2; targeted static inspection.\par
- Replay records: , and ; and . Re-run only reviewer-authored document-processing scripts if desired.\par
Initial-review and proposal build history. An independent AI checker reviewed the candidate issue and alternative interpretation. This is not human validation or a formal publication decision. Human scientific adjudication is pending; no author response, correction acceptance or adoption is claimed. The PDF was rendered from the same canonical content as the self-contained TeX using an existing PDF renderer. No TeX compiler was available locally; controller compilation was pending at that stage.\par
Actual replay commands, run from verification/0028, are shown below. The document checks completed successfully and confirmed the historical paraphrase experiment on PDF pp. 8–9. Original-source hashes, supporting-source hashes and retrieved snapshots are recorded in , and acquisition manifests. Public metadata retrieval is time-dependent.\par
> python -X utf8 python -X utf8\par
Initial-review and proposal build history. The direct PDF reader had failed earlier in this task; existing Poppler extraction and page rendering succeeded. OpenReview remained behind browser verification. These access limits, the lack of local TeX compilation, and the absence of human adjudication are preserved in the evidence records.\par
Batch token: [internal identifier omitted]. Scientific decisions and human verification remain pending; author adoption was not requested.\par
Controller build revision. The current report PDF was compiled from this self-contained TeX using the existing Tectonic compiler in untrusted mode. Long URLs, hashes and file paths were made breakable; scientific content is unchanged by these layout repairs. Earlier previews and build records remain preserved. Current build hashes and verification results are in the accompanying verification files. Human scientific adjudication remains pending.\par
\end{document}
